\chapter{Uniform Integrability and Conditional Limits}
\label{sec:p15-uniform-integrability-and-conditional-limits}

Almost-sure convergence alone does not justify passing a limit through
expectation. Uniform integrability supplies exactly the missing control
when the variables already converge in probability. It also explains
when growing or decreasing information determines a conditional limit.
The source presentation developed here is Durrett's Section~4.6 and
the first three theorems of Section~4.7
\cite[pp.~244--250]{DurrettPTEAuthor2019}.

\section{Uniform control of integrable tails}
\label{sec:p15-ui-definitions}

\begin{hbdefinition}{def:p15-uniform-integrability}{Uniform integrability and small-event control}
Work on a probability space \((\Omega,\mathcal F,P)\).
A family \((X_i)_{i\in I}\) of real integrable random variables is
\emph{uniformly integrable} if
\[
       \lim_{M\to\infty}\sup_{i\in I}
                 E\bigl(|X_i|\mathbf1_{\{|X_i|>M\}}\bigr)=0.
\]
For an empty family each displayed supremum of nonnegative quantities
is defined to be zero. The family has \emph{uniform small-event control} if
for every \(\epsilon>0\) there is a \(\delta>0\) such that
\[
  P(A)\leq\delta\quad\Longrightarrow\quad
  \sup_{i\in I}E(|X_i|\mathbf1_A)\leq\epsilon
                 \qquad(A\in\mathcal F).
\]
Write \(E(Z;A)=E(Z\mathbf1_A)\) when this integral is defined.
All conditional expectations below are versions equal almost surely;
no completeness of a sub-sigma-algebra is assumed.
\index{uniform integrability}\index{conditional expectation!limits}
\end{hbdefinition}

\begin{hbtheorem}{thm:p15-ui-small-sets}{proved=annex:p15.ui-small-sets}{Equivalent tail and small-event conditions}
A family is uniformly integrable if and only if
\[
          \sup_i E|X_i|<\infty
       \quad\hbox{and it has uniform small-event control.}
\]
Every family dominated in absolute value by one integrable variable is
uniformly integrable. Every finite family of integrable variables is
uniformly integrable. Bounded first moments alone do not imply this property.
\end{hbtheorem}

\begin{hbtheorem}{thm:p15-ui-moment-criteria}{proved=annex:p15.ui-moment-criteria}{Superlinear moment bounds and their converse}
Let \(\Phi:[0,\infty)\to[0,\infty)\) be finite and Borel measurable,
with \(\Phi(t)/t\to\infty\) as \(t\to\infty\).
If \(\sup_i E\Phi(|X_i|)<\infty\), the family is uniformly integrable.
Examples are \(\Phi(t)=t^p\), \(p>1\), and
\(\Phi(t)=t\log^+t\), where \(\log^+t=\max(0,\log t)\) for \(t>0\)
and \(\Phi(0)=0\).

Conversely, for every uniformly integrable family there exists such a
\(\Phi\) that is continuous, convex and nondecreasing, has \(\Phi(0)=0\),
and satisfies \(\sup_i E\Phi(|X_i|)\leq1\).
The converse is supplied by an explicit sum of positive-part functions.
\index{de la Vall\'ee Poussin criterion}
\end{hbtheorem}

\begin{hbtheorem}{thm:p15-ui-conditional-family}{proved=annex:p15.ui-conditional-family}{All conditional expectations of one integrable variable}
For \(X\in L^1(P)\), the family
\[
 \{E(X\mid\mathcal G):\mathcal G
                 \text{ is a sub-sigma-algebra of }\mathcal F\}
\]
is uniformly integrable. In particular this conclusion needs no nesting
or countability assumption on the family of information sets
\cite[thm.~4.6.1]{DurrettPTEAuthor2019}.
\end{hbtheorem}

\section{The exact criterion for convergence in mean}
\label{sec:p15-ui-convergence}

\begin{hbtheorem}{thm:p15-probability-l1-tools}{proved=annex:p15.probability-l1-tools}{Subsequences, bounded limits and integrals over events}
For real random variables the following implications hold.
\begin{enumerate}
\item Convergence in \(L^1\) implies convergence in probability.
\item From a sequence converging in probability one can select a
      subsequence converging almost surely to the same limit.
\item If \(Z_n\to Z\) in probability and \(|Z_n|,|Z|\leq K<\infty\),
      then \(E|Z_n-Z|\to0\).
\item If \(X_n\to X\) in probability, then
      \(E|X|\leq\liminf_n E|X_n|\).
\item If \(X_n\to X\) in \(L^1\), then, uniformly over \(A\in\mathcal F\),
      \[
        |E(X_n;A)-E(X;A)|\leq E|X_n-X|\longrightarrow0.
      \]
\end{enumerate}
The fourth assertion permits infinite expectations.
\end{hbtheorem}

\begin{hbtheorem}{thm:p15-vitali-convergence}{proved=annex:p15.vitali-convergence}{Uniform integrability, mean convergence and absolute means}
Suppose \(X_n\in L^1\) and \(X_n\to X\) in probability, with \(X\)
real and measurable. The following are equivalent:
\begin{enumerate}
\item \((X_n)_{n\geq0}\) is uniformly integrable;
\item \(X\in L^1\) and \(E|X_n-X|\to0\);
\item \(E|X_n|\to E|X|<\infty\).
\end{enumerate}
Both source truncations are retained: clipping to \([-M,M]\) proves
the first implication, and the continuous cutoff of the absolute value
proves the third-to-first implication
\cite[thm.~4.6.3]{DurrettPTEAuthor2019}.
\index{Vitali convergence theorem}
\end{hbtheorem}

\begin{hbtheorem}{thm:p15-submartingale-ui}{proved=annex:p15.submartingale-ui}{Convergence in mean for submartingales}
For a submartingale, uniform integrability, convergence almost surely
and in \(L^1\), and convergence in \(L^1\) are equivalent.
The same equivalence holds for supermartingales by changing sign.
In the uniformly integrable case the real integrable limit \(X_\infty\)
may be chosen measurable with respect to
\(\mathcal F_\infty=\sigma(\bigcup_{n\geq0}\mathcal F_n)\), and
\[
        X_n\leq E(X_\infty\mid\mathcal F_n)
\]
for submartingales, with the reverse inequality for supermartingales.
\end{hbtheorem}

\begin{hbdefinition}{def:p15-closed-martingale}{An integrable closing variable}
A martingale \((X_n,\mathcal F_n)\) is \emph{closed} by \(Z\in L^1(P)\)
if \(X_n=E(Z\mid\mathcal F_n)\) for every \(n\).
Write \(\mathcal F_\infty=\sigma(\bigcup_n\mathcal F_n)\).
The variable \(Z\) need not be measurable with respect to
\(\mathcal F_\infty\).
\index{martingale!closed}
\end{hbdefinition}

\begin{hbtheorem}{thm:p15-closed-martingale}{proved=annex:p15.closed-martingale}{The four equivalent martingale conditions}
For a martingale the following are equivalent: uniform integrability;
convergence almost surely and in \(L^1\); convergence in \(L^1\);
and existence of an integrable closing variable.
If \(X_n\to X_\infty\) in \(L^1\), then
\[
                       X_n=E(X_\infty\mid\mathcal F_n).
\]
If \(Z\) closes the martingale, then
\(X_\infty=E(Z\mid\mathcal F_\infty)\) almost surely
\cite[lems.~4.6.5--6, thm.~4.6.7]{DurrettPTEAuthor2019}.
\end{hbtheorem}

\section{Increasing information and its applications}
\label{sec:p15-increasing-information}

\begin{hbtheorem}{thm:p15-levy-upward}{proved=annex:p15.levy-upward}{L\'evy's upward conditional-limit theorem}
If \(Z\in L^1\), \(\mathcal F_n\) increases, and
\(\mathcal F_\infty=\sigma(\bigcup_n\mathcal F_n)\), then
\[
       E(Z\mid\mathcal F_n)\longrightarrow
       E(Z\mid\mathcal F_\infty)
                  \quad\hbox{almost surely and in }L^1.
\]
The limit is defined measurably even on the exceptional set where
the initial versions fail to converge
\cite[thm.~4.6.8]{DurrettPTEAuthor2019}.
\end{hbtheorem}

\begin{hbtheorem}{thm:p15-levy-information-law}{proved=annex:p15.levy-information-law}{Learning a measurable event and the tail zero-one law}
Let \(\mathcal F_n\) increase and
\(\mathcal F_\infty=\sigma(\bigcup_n\mathcal F_n)\).
If \(A\in\mathcal F_\infty\), then
\[
                P(A\mid\mathcal F_n)=E(\mathbf1_A\mid\mathcal F_n)
                         \longrightarrow\mathbf1_A
\]
almost surely and in \(L^1\). Consequently, for independent random
variables \(\xi_i\), every event in
\(\bigcap_m\sigma(\xi_m,\xi_{m+1},\ldots)\) has probability zero or one.
This supplies the source's separate conditional-limit proof of
Kolmogorov's law
\cite[thm.~4.6.9 and continuation]{DurrettPTEAuthor2019}.
\index{L\'evy zero-one law}
\end{hbtheorem}

\begin{hbtheorem}{thm:p15-changing-conditional-limits}{proved=annex:p15.changing-conditional-limits}{Changing the variable and the information together}
Let \(\mathcal F_n\) increase and
\(\mathcal F_\infty=\sigma(\bigcup_n\mathcal F_n)\).
\begin{enumerate}
\item If \(Y_n\to Y\) in \(L^1\), then
      \(E(Y_n\mid\mathcal F_n)\to E(Y\mid\mathcal F_\infty)\) in \(L^1\).
\item If \(Y_n\to Y\) almost surely and \(|Y_n|\leq Z\) for all \(n\),
      with \(Z\geq0\) integrable, then the same conditional sequence
      converges almost surely and in \(L^1\).
\end{enumerate}
In the second assertion \(|Y|\leq Z\) almost surely.
Its almost-sure proof uses the decreasing oscillation envelopes
\[
             W_N=\sup_{n,m\geq N}|Y_n-Y_m|,
\]
not an interchange of an uncountable family of conditional null sets
\cite[thm.~4.6.10, ex.~4.6.7]{DurrettPTEAuthor2019}.
\end{hbtheorem}

\begin{hbtheorem}{thm:p15-ui-counterexamples}{proved=annex:p15.ui-counterexamples}{Three limits of the hypotheses}
\begin{enumerate}
\item On \((0,1)\) with uniform probability,
      \(X_n=n\mathbf1_{(0,1/n)}\) tends to zero almost surely and
      has \(E|X_n|=1\), but is not uniformly integrable.
\item A common fair sign \(\varepsilon\) gives
      \(X_n=(-1)^n\varepsilon\). The family is uniformly integrable
      and all absolute means equal one, but it has no limit in probability.
\item Construct mutually independent Bernoulli variables \(B_n,C_n\),
      \(n\geq1\), with
      \(P(B_n=1)=P(C_n=1)=1/n\). Let \(X_n=nB_nC_n\) and
      \(\mathcal G=\sigma(B_1,B_2,\ldots)\). Then \(X_n\to0\)
      almost surely and in \(L^1\), and \((X_n)\) is uniformly integrable.
      Nevertheless
      \[
                      E(X_n\mid\mathcal G)=B_n
      \]
      fails to converge almost surely, although it converges to zero in
      \(L^1\).
\end{enumerate}
The third example is the complete construction of source Example~4.6.11.
\end{hbtheorem}

\begin{hbdefinition}{def:p15-dyadic-information}{Dyadic information on the unit interval}
On \(\Omega=[0,1)\) with its Borel sets and uniform probability, set
\[
 I_{k,n}=[k2^{-n},(k+1)2^{-n}),\qquad
 \mathcal D_n=\sigma(I_{k,n}:0\leq k<2^n),\qquad n\geq0.
\]
Each atom has mass \(2^{-n}\). Its right endpoint may equal one.
\end{hbdefinition}

\begin{hbtheorem}{thm:p15-dyadic-martingales}{proved=annex:p15.dyadic-martingales}{Dyadic averages, step approximation and Lipschitz increments}
For Borel \(f\in L^1([0,1))\), the conditional average
\[
       E(f\mid\mathcal D_n)(x)=
                2^n\int_{I_{k,n}}f(t)\,dt\quad(x\in I_{k,n})
\]
converges to \(f\) almost everywhere and in \(L^1\).
Thus dyadic step functions approximate every such \(f\) in \(L^1\).
The conclusion also holds for functions measurable in the completed
Lebesgue sigma-algebra, using their Borel versions.

If \(f:[0,1)\to\mathbb R\) is Lipschitz with constant \(K\geq0\), meaning
\(|f(s)-f(t)|\leq K|s-t|\) for all \(s,t\in[0,1)\), extend it
continuously to \(1\). The slopes
\[
       X_n(x)=2^n\bigl[f((k+1)2^{-n})-f(k2^{-n})\bigr],
                         \quad x\in I_{k,n},
\]
form a bounded martingale. Its limit \(X_\infty\) satisfies
\(|X_\infty|\leq K\), convergence almost everywhere and in \(L^1\), and
\[
             f(b)-f(a)=\int_a^b X_\infty(t)\,dt
                                  \quad(0\leq a\leq b\leq1).
\]
These are the complete source Exercises~4.6.2--3, including the endpoint.
\end{hbtheorem}

\begin{hbtheorem}{thm:p15-posterior-mean-consistency}{proved=annex:p15.posterior-mean-consistency}{Recovery of a common additive parameter}
Let \(\xi_i\) be i.i.d. with \(E|\xi_i|<\infty\) and known mean \(\mu\).
Let \(\theta\in L^1\) be independent of the sequence, put
\(Y_i=\theta+\xi_i\), and set
\(\mathcal F_n=\sigma(Y_1,\ldots,Y_n)\).
Then
\[
                    E(\theta\mid\mathcal F_n)\longrightarrow\theta
                           \quad\hbox{almost surely and in }L^1.
\]
In particular this applies to independent normal \((0,1)\) noise.
It is a theorem for the stated additive sampling model and integrable
prior variable, rather than an unrestricted posterior consistency claim
\cite[ex.~4.6.1]{DurrettPTEAuthor2019}.
\end{hbtheorem}

\begin{hbtheorem}{thm:p15-extinction-or-escape}{proved=annex:p15.extinction-or-escape}{A positive chance of absorption prevents bounded survival}
Let \(X_n\) be finite nonnegative random variables, \(n\geq1\),
\(\mathcal F_n=\sigma(X_1,\ldots,X_n)\), and
\(D=\bigcup_{n\geq1}\{X_n=0\}\).
Suppose for every positive integer \(K\) there is \(\delta_K>0\) with
\[
       P(D\mid\mathcal F_n)\geq\delta_K
                              \quad\hbox{on }\{X_n\leq K\}
                              \quad\text{for every }n.
\]
Then \(P(D\cup\{X_n\to\infty\})=1\).

For a finite-population Galton--Watson process, define
\(Z_{n+1}=\sum_{j=1}^{Z_n}\xi_{n,j}\), where the offspring variables
are i.i.d. nonnegative integers, independent of an almost-surely finite
nonnegative integer \(Z_0\),
and let \(p_0=P(\xi_{n,j}=0)>0\).
No offspring mean assumption is needed to conclude
\[
             P(Z_n\text{ is eventually zero, or }Z_n\to\infty)=1.
\]
The empty offspring sum is zero
\cite[exs.~4.6.4--5]{DurrettPTEAuthor2019}.
\end{hbtheorem}

\begin{hbtheorem}{thm:p15-bounded-polarization}{proved=annex:p15.bounded-polarization}{A bounded two-outcome martingale}
Let \(0<\alpha,\beta<1\), \(\alpha+\beta=1\), and let
\(X_n\in[0,1]\) be adapted with
\[
 \begin{split}
 P(X_{n+1}=\alpha+\beta X_n\mid\mathcal F_n)&=X_n,\\
 P(X_{n+1}=\beta X_n\mid\mathcal F_n)&=1-X_n.
 \end{split}
\]
Then \(X_n\) converges almost surely and in \(L^1\) to a variable
in \(\{0,1\}\), and \(P(X_\infty=1)=EX_0\).
In particular, for deterministic \(X_0=\theta\in[0,1]\), the probability
is \(\theta\) \cite[ex.~4.6.6]{DurrettPTEAuthor2019}.
\end{hbtheorem}

\section{Decreasing information}
\label{sec:p15-decreasing-information}

\begin{hbdefinition}{def:p15-backward-martingale}{Reverse time and its index conventions}
Let \(\mathcal G_0\supseteq\mathcal G_1\supseteq\cdots\)
be sub-sigma-algebras of \(\mathcal F\).
A \emph{backward martingale} in nonnegative reverse-time indices is a
sequence of integrable \(\mathcal G_n\)-measurable variables \(Y_n\) with
\[
                        E(Y_n\mid\mathcal G_{n+1})=Y_{n+1}.
\]
For \(m\geq n\) the tower identity then gives
\(E(Y_n\mid\mathcal G_m)=Y_m\).
Set \(\mathcal G_\infty=\bigcap_n\mathcal G_n\).
In the source's negative-index notation \(X_{-n}=Y_n\) and
\(\mathcal F_{-n}=\mathcal G_n\); consequently
\(E(X_{k+1}\mid\mathcal F_k)=X_k\) for \(k\leq-1\).
\index{martingale!backward}\index{martingale!reverse}
\end{hbdefinition}

\begin{hbtheorem}{thm:p15-backward-martingale-convergence}{proved=annex:p15.backward-martingale-convergence}{Every backward martingale converges in mean}
A backward martingale is uniformly integrable and has a real
\(\mathcal G_\infty\)-measurable version of its almost-sure and \(L^1\)
limit. That version satisfies
\[
                    Y_\infty=E(Y_0\mid\mathcal G_\infty).
\]
The proof applies the forward upcrossing inequality to each finite
string \(Y_n,Y_{n-1},\ldots,Y_0\), with its corresponding increasing
filtration \cite[thms.~4.7.1--2]{DurrettPTEAuthor2019}.
\end{hbtheorem}

\begin{hbtheorem}{thm:p15-levy-backward}{proved=annex:p15.levy-backward}{L\'evy's downward conditional-limit theorem}
For \(Z\in L^1\) and decreasing \(\mathcal G_n\),
\[
              E(Z\mid\mathcal G_n)\longrightarrow
              E(Z\mid\mathcal G_\infty)
                        \quad\hbox{almost surely and in }L^1.
\]
All preceding upward and downward conditional limits apply coordinate
by coordinate to a finite-dimensional integrable random vector.
No interchange of an uncountable coordinate family is asserted
\cite[thm.~4.7.3]{DurrettPTEAuthor2019}.
\end{hbtheorem}

The complete Section~4.7 applications and exercises are supplied in
\cref{sec:p15-exchangeability-and-definetti}, including the conditional
product theorem, the standard Borel mixture, the ballot identity and
both reverse-limit exercises. The present reverse-time theorem
supplies their convergence premise. Its application to Blackwell's general
comparison of experiments is completed in
\cref{thm:p17-blackwell-joint-process,thm:p17-general-convex-transport}.
